\chapter{Flujo de competencia}
\label{chap:flujo-competencia}

\section{Vision completa del ciclo}
\label{sec:flujo-vision}

El flujo de competencia empieza antes del panel de torneo. La plataforma construye el torneo a partir de registros publicos confirmados. Un equipo solo entra al dominio competitivo despues de confirmar asistencia y sincronizarse como \clase{Team}. Desde ese punto, los servicios de torneo crean grupos, batallas, estados, historial y podio.

El ciclo completo es:

\begin{enumerate}
  \item inscripcion publica;
  \item validacion de formulario;
  \item confirmacion de asistencia por token;
  \item creacion o actualizacion de equipo operativo;
  \item seleccion de division;
  \item inicializacion de competencia;
  \item fase de grupos;
  \item clasificacion;
  \item eliminatorias, repechajes o fases progresivas;
  \item correcciones manuales cuando corresponden;
  \item sincronizacion de estados e historial;
  \item final y podio.
\end{enumerate}

\begin{figure}[H]
  \centering
  \fbox{\parbox{0.88\textwidth}{\centering Marcador de diagrama: flujo completo del torneo desde registro hasta ganador final. Fuente prevista: DIA-08 en \archivo{planificacion/06_inventario_diagramas.md}.}}
  \caption{Marcador para flujo completo de competencia.}
  \label{fig:flujo-completo-torneo}
\end{figure}

\section{Inscripcion, validacion y asistencia}
\label{sec:flujo-inscripcion}

El flujo inicia en \archivo{apps.participants}. \clase{SchoolRegistrationForm} y \clase{UniversityRegistrationForm} validan que exista edicion activa, que los integrantes tengan datos consistentes, que el nombre del robot no duplique variantes normalizadas en la edicion, que el correo no se repita y que los documentos no existan ya en otros registros de la edicion.

Al guardar, se crea \clase{SchoolRegistration} o \clase{UniversityRegistration} con estado \variable{submitted}, junto con sus participantes. Ese registro aun no es un equipo de competencia. El cambio ocurre en la confirmacion de asistencia: \funcion{confirm_attendance} marca asistencia, aprueba registro, llama \funcion{sync_registration_to_team} y puede invocar \funcion{initialize_competition} para la division correspondiente.

La razon de este diseño es evitar competencias con inscritos que no asistieron. El dato que alimenta el torneo no es la intencion de participar, sino la asistencia confirmada y sincronizada.

\section{Divisiones y competencia operativa}
\label{sec:flujo-divisiones}

El sistema separa competencias por division mediante \clase{DivisionType}: colegios y universidades. Cada \clase{DivisionCompetition} pertenece a una edicion y una division, con restriccion de unicidad por par edicion/division. Esto evita que el panel tenga dos competencias operativas simultaneas para la misma categoria.

La competencia contiene:

\begin{itemize}
  \item \variable{format_key}: perfil aplicado;
  \item \variable{configuration}: configuracion flexible de perfil, stages, propuestas, podio y otros datos dinamicos;
  \item \variable{status}: estado global de competencia;
  \item \variable{shuffle_seed}: semilla usada para asignaciones.
\end{itemize}

Los equipos entran a la competencia desde \clase{Team}. Los servicios crean grupos, entries, batallas y estados por equipo. La persistencia conserva tanto estructura como avance.

\section{Inicializacion y perfiles}
\label{sec:flujo-inicializacion}

\funcion{initialize_competition} obtiene equipos confirmados de la edicion y division, calcula perfil, crea o reutiliza \clase{DivisionCompetition} y llama a servicios de reconstruccion si es seguro hacerlo. Si la competencia ya tiene progreso o historial, el servicio bloquea cambios de configuracion que podrian invalidar resultados.

\begin{longtable}{p{3.2cm}p{9.5cm}}
\toprule
Cantidad de equipos & Estructura aplicada \\
\midrule
Menos de 4 & Modo pendiente; no se genera torneo automatico. \\
4 a 15 & 2 grupos, 2 clasificados por grupo, semifinal y final. \\
16 a 31 & 4 grupos, 2 clasificados, cuartos, semifinal y final. \\
32 a 63 & 8 grupos, 2 clasificados y repechaje 1. \\
64 o mas & 16 grupos, ronda de 32 y repechaje 1. \\
\bottomrule
\end{longtable}

La configuracion manual del perfil permite ajustar grupos y clasificados desde el panel. Sin embargo, esos ajustes quedan sometidos a validaciones de progreso. La plataforma privilegia consistencia de resultados sobre flexibilidad tardia.

\section{Fase de grupos}
\label{sec:flujo-grupos}

La fase de grupos usa \clase{DivisionGroup} y \clase{DivisionGroupEntry}. La asignacion balanceada procura distribuir equipos y separar instituciones. En el panel, el operador selecciona clasificados por grupo. \funcion{set_group_qualifiers} exige que la cantidad seleccionada coincida con \variable{qualifiers_per_group}.

Cuando se guardan clasificados:

\begin{itemize}
  \item se actualizan entries del grupo;
  \item se registra historial por equipo;
  \item se sincronizan estados de competencia;
  \item se habilita la construccion de fases posteriores si los grupos estan completos.
\end{itemize}

Si ya hay progreso posterior, el servicio puede exigir \variable{manual_override}. Esta proteccion evita que una correccion de grupos borre implicitamente el significado de batallas ya creadas.

\begin{figure}[H]
  \centering
  \fbox{\parbox{0.88\textwidth}{\centering Marcador de diagrama: transicion entre fase de grupos y fases posteriores. Fuente prevista: DIA-09 y DIA-10.}}
  \caption{Marcador para transicion entre fases.}
  \label{fig:flujo-transicion-fases}
\end{figure}

\section{Eliminatorias}
\label{sec:flujo-eliminatorias}

Las eliminatorias se basan en \clase{CompetitionBattle}. Una batalla puede ser duelo, triangular, campal o grupo. Las etapas se identifican con \clase{CompetitionStage}: ronda de 32, octavos, cuartos, semifinal, tercer lugar y final, entre otras.

El avance se materializa con servicios de siembra y recomendacion. \funcion{seed_duel_stage} carga cruces desde clasificados. \funcion{seed_next_duel_stage} usa ganadores de fase anterior. \funcion{wire_final_stages} conecta semifinal, tercer lugar y final cuando aplica. Los resultados se guardan con \funcion{set_battle_winner} o \funcion{set_battle_qualifiers}.

La informacion evoluciona en tres niveles:

\begin{enumerate}
  \item batalla: ganador o clasificados;
  \item estado de equipo: stage y status actual;
  \item historial: evento que explica el movimiento.
\end{enumerate}

Este triple registro permite mostrar el estado actual sin perder la explicacion del avance.

\section{Repechajes}
\label{sec:flujo-repechajes}

El repechaje permite reintegrar participantes bajo condiciones controladas. \funcion{materialize_repechage_stage} sincroniza estados, calcula candidatos, define formato seguro y revisa si ya existe una etapa abierta. Si ya hay un repechaje disponible, el servicio evita crear duplicados.

El flujo de repechaje no es un formulario independiente: depende de estados eliminados, recientes o repechables. Por eso se apoya en \archivo{apps/tournament/recommendations.py}, que calcula preview y alternativas. La materializacion crea batallas y marca estados como \variable{repechage} cuando corresponde.

\section{Fases manuales y correcciones}
\label{sec:flujo-fases-manuales}

Las fases manuales cubren escenarios donde la recomendacion automatica no debe decidir sola. El operador genera una propuesta, revisa distribucion, puede reasignar grupos y confirma. La propuesta se guarda temporalmente en \variable{DivisionCompetition.configuration}.

El servicio valida firmas antes de materializar:

\begin{itemize}
  \item firma de propuesta;
  \item firma de participantes elegibles;
  \item etapa origen;
  \item grupos no vacios;
  \item participantes sin duplicados;
  \item clasificados validos por grupo.
\end{itemize}

La correccion manual de resultados usa otro mecanismo: \clase{ManualCorrectionRequired}. Si una accion puede afectar progreso posterior, el servicio detiene la operacion y la vista responde con \variable{requires_confirmation}. \archivo{control.js} muestra confirmacion y reintenta con \variable{manual_override}. Esto conserva control humano sobre cambios riesgosos sin eliminar validaciones de consistencia.

\begin{figure}[H]
  \centering
  \fbox{\parbox{0.88\textwidth}{\centering Marcador de diagrama: correccion manual con excepcion controlada, confirmacion AJAX y reintento con \variable{manual_override}. Fuente prevista: DIA-12 y DIA-13.}}
  \caption{Marcador para correccion manual y AJAX.}
  \label{fig:flujo-correccion-manual}
\end{figure}

\section{Actualizacion de estados}
\label{sec:flujo-estados}

\funcion{sync_team_states} y \funcion{sync_competition} mantienen coherencia entre estructura y estado. La competencia no depende solo de saber que batalla existe; necesita saber donde esta cada equipo y que estado tiene: activo, eliminado, repechaje, clasificado o retirado.

\clase{TeamCompetitionState} actua como fotografia actual. \clase{CompetitionHistoryEntry} conserva eventos. Cuando se guarda un resultado, la plataforma actualiza ambas capas. Esta separacion permite que el panel muestre participantes por etapa y que el modal de participante muestre historial.

\begin{figure}[H]
  \centering
  \fbox{\parbox{0.88\textwidth}{\centering Marcador de diagrama: propagacion de estados desde grupos y batallas hacia TeamCompetitionState e historial.}}
  \caption{Marcador para propagacion de estados.}
  \label{fig:flujo-propagacion-estados}
\end{figure}

\section{Historial de competencia}
\label{sec:flujo-historial}

Cada movimiento relevante puede generar \clase{CompetitionHistoryEntry}. El historial registra accion, stage, status, titulo, descripcion, estados previos y nuevos, grupos, batalla y metadata. Esto permite reconstruir decisiones sin depender de memoria del operador.

La conservacion de historial tambien protege configuraciones posteriores. El codigo revisa si hay historial para decidir si una competencia puede reconstruirse o si un perfil puede cambiar. Por tanto, el historial no solo sirve para auditoria visual; participa en decisiones de consistencia.

\begin{figure}[H]
  \centering
  \fbox{\parbox{0.88\textwidth}{\centering Marcador de diagrama: persistencia del historial de competencia y uso en bloqueos de reconstruccion.}}
  \caption{Marcador para persistencia del historial.}
  \label{fig:flujo-persistencia-historial}
\end{figure}

\section{Cierre y ganador final}
\label{sec:flujo-cierre}

La competencia se cierra con \funcion{save_final_podium}. El servicio valida que exista final, que las posiciones pertenezcan a participantes de la final, que el tercer lugar exista cuando la final tenga tres o mas participantes y que las posiciones sean equipos distintos. Luego guarda \variable{final_podium} en \variable{competition.configuration}, marca la competencia como \variable{completed} y registra historial para las posiciones.

El ganador final no es solo un campo aislado. Queda conectado a la batalla final, al podio persistido, al estado de competencia y al historial. Esa redundancia controlada facilita mostrar el resultado final y revisar como se obtuvo.

\section{Escenarios alternos relevantes}
\label{sec:flujo-escenarios-alternos}

La auditoria de calidad identifico que el manual debe cubrir escenarios no lineales. En este bloque se documentan los mecanismos que los soportan:

\begin{itemize}
  \item \textbf{Tres finalistas}: \funcion{save_final_podium} exige tercer lugar cuando la final tiene tres o mas participantes.
  \item \textbf{Multiples clasificados}: \funcion{set_battle_qualifiers} maneja batallas donde clasifican varios equipos.
  \item \textbf{Repechaje existente}: \funcion{materialize_repechage_stage} revisa fases abiertas y evita duplicados.
  \item \textbf{Fase manual cancelada}: \funcion{cancel_manual_phase_proposal} elimina una propuesta pendiente sin materializarla.
  \item \textbf{Participantes cambiantes}: las firmas de propuesta manual fuerzan regenerar si los elegibles cambiaron.
  \item \textbf{Correccion posterior}: \clase{ManualCorrectionRequired} exige confirmacion antes de modificar resultados con progreso posterior.
\end{itemize}

\section{Impacto de modificacion}
\label{sec:flujo-impacto-modificacion}

\begin{longtable}{p{3.7cm}p{4.2cm}p{2.6cm}p{3.4cm}}
\toprule
Componente & Componentes relacionados & Riesgo & Accion recomendada \\
\midrule
Registro y asistencia & Participantes, equipos, inicializacion & Alto & Probar confirmacion y sincronizacion por division. \\
Perfiles automaticos & Inicializacion, grupos, stages, vistas & Alto & Probar rangos de equipos y overrides manuales. \\
Fase de grupos & Entries, estados, historial, siguiente fase & Alto & Probar cantidad exacta de clasificados y correccion. \\
Eliminatorias & Batallas, ganadores, podio, historial & Alto & Probar duelos, multiples clasificados y fases posteriores. \\
Repechaje & Estados eliminados, recomendaciones, batallas & Medio & Probar candidatos 2, 3, 4, 6 y fase ya existente. \\
Fase manual & Configuration, firmas, elegibles, batallas & Alto & Probar generacion, reasignacion, confirmacion y cancelacion. \\
Historial & Estados, modales, reconstruccion & Medio & Verificar que cada accion critica registre evento. \\
Podio & Final, configuration, historial, status & Alto & Probar posiciones distintas y final con tercer lugar. \\
\bottomrule
\end{longtable}

\section{Consideraciones de mantenimiento}
\label{sec:flujo-mantenimiento}

\begin{longtable}{p{4cm}p{9cm}}
\toprule
Aspecto & Consideracion \\
\midrule
Archivos relacionados & \archivo{apps/participants/services.py}, \archivo{apps/tournament/services.py}, \archivo{formats.py}, \archivo{recommendations.py}, \archivo{views.py}, \archivo{models.py}, \archivo{static/js/control.js} y templates de torneo. \\
Pruebas recomendadas & Registro, asistencia, inicializacion, perfiles, grupos, eliminatorias, repechaje, fase manual, correccion AJAX, propagacion de estados y podio. \\
Riesgos & Cambios en una etapa pueden invalidar estados posteriores, historial, configuration o UI del panel. \\
Dependencias & \clase{Team}, \clase{DivisionCompetition}, \clase{DivisionGroup}, \clase{CompetitionBattle}, \clase{TeamCompetitionState}, \clase{CompetitionHistoryEntry} y \archivo{control.js}. \\
Buenas practicas & No modificar flujo por pantalla aislada; revisar servicios, modelos, vistas, templates y JavaScript como una unidad; mantener historial y validar cambios con datos de ejemplo. \\
\bottomrule
\end{longtable}
